\ptitle{Principalele mulțimi de numere}
    \psubtitle{Mulțimea numerelor naturale ($\mathbb{N}$)}
        Mulțimea numerelor naturale este formate din numerele mai mare sau egale cu 0 și fără o parte decimală.
        \[ \mathbb{N} = \{0, 1, 2, 3, ..., \infty\} \]
        Definim și mulțimea numerelor naturale nenule ($\mathbb{N}^*$) în următorul fel.
        \[ \mathbb{N}^* = \mathbb{N} / \{0\} = \{1, 2, 3, ..., \infty\} \]
    \psubtitle{Mulțimea numerelor întregi ($\mathbb{Z}$)}
        Mulțimea numerelor întregi este formată din toate numerele fără o parte decimală sau prin reuniunea dintre mulțimea numerelor naturale cu opusul ei.
        \[ \mathbb{Z} = \{ -\infty, ..., -2, -1, 0, 1, 2, ..., \infty \} = \mathbb{N} \cup -\mathbb{N}\]
        Mai definim și mulțimea numerelor întregi negative ($\mathbb{Z}_-$) și pozitive ($\mathbb{Z}_+$).
        \[ \mathbb{Z}_- = \mathbb{Z} / \mathbb{N}, \qquad \mathbb{Z}_+ = \mathbb{N}^* \]
        Reuniune celor doua rezultă în mulțimea numerelor întregi nenule ($\mathbb{Z}^*$).
        \[ \mathbb{Z}^* = \mathbb{Z}_- \cup \mathbb{Z}_+ = \mathbb{Z} / \{0\} \] 
    \psubtitle{Mulțimea numerelor raționale ($\mathbb{Q}$)}
        Mulțimea numerelor raționale este formată din toate numerele care se pot forma prin împărțirea a două numere întregi.
        \[ \mathbb{Q} = \left\{ \left. \frac{a}{b} \right| a,b \in \mathbb{Z}, b \neq 0 \right\} \]
        Definim și mulțimea numerelor raționale negative ($\mathbb{Q}_-$) și pozitive ($\mathbb{Q}_+$).
        \[ \mathbb{Q}_- = \left\{ \left. -\frac{a}{b}\right| a, b \in \mathbb{N}^* \right\}, \qquad 
        \mathbb{Q}_+ = \mathbb{Q} / \mathbb{Q}_- / \{0\} \]
        Reuniune celor două rezultă în mulțimea numerelor raționale nenule ($\mathbb{Q}^*$).
        \[ \mathbb{Q}^* = \mathbb{Q}_- \cup \mathbb{Q}_+ = \mathbb{Q} / \{0\} \]
    \psubtitle{Mulțimea numerelor iraționale ($\mathbb{R} / \mathbb{Q}$)}
        Mulțimea numerelor iraționale este formată din toate numerele care nu pot fii scrise ca a fracție dintre două numere întregi. \\
        Exemple de numere iraționale.
        \[ \{ \pm \sqrt{n}, \pi, \phi \} \subset \mathbb{R} / \mathbb{Q}, n \not - \text{pp.}\] 
        Exemple de numere \textbf{raționale}, dar unele presoane cred că ele sunt iraționale.
        \[ \{ 0,(9), 0,(3) \} \subset \mathbb{Q} \not \subset \mathbb{R} \]
    \psubtitle{Mulțimea numerelor reale ($\mathbb{R}$)}
        Mulțimea numerelor reale este cea mai mare mulțimea învățata, ea este formată prin reuniunea mulțimii numerelor iraționale ($\mathbb{R} / \mathbb{Q}$) cu mulțimea numerelor raționale ($\mathbb{Q}$).
        \[ \mathbb{R} = (\mathbb{R} / \mathbb{Q}) \cup \mathbb{Q} \]
        Definim mulțimile numerelor reale pozitive ($\mathbb{R}_+$) și negarive ($\mathbb{R}_-$) în felul următor.
        \[ \mathbb{R}_+ = \{ x | x \in \mathbb{R}, x > 0 \}, \qquad \mathbb{R}_- = \mathbb{R} / \mathbb{R}_+ / \{0\} \]
    \psubtitle{Relația de incluziune dintre principalele mulțimi de numere}
        Mulțiimile principale sunt incluse una în următoarea în ordine în care le-am definit (excluzând mulțimea numerelor iraționale, care este inclusa doar în mulțimea numerelor reale).
        \[ \mathbb{N} \subset \mathbb{Z} \subset \mathbb{Q} \subset \mathbb{R}, \quad (\mathbb{R} / \mathbb{Q}) \subset \mathbb{R} \]
        \textbf{Notă}: Relația de mai sus nu înseamnă că mulțimile care sunt incluse în alte mulțimi sunt mai mici decât ele, deoarece mulțimile date nu sunt finite, ci infinite.
